% LaTeX source maintained directly for Overleaf.
\chapter{IV. Truy vấn, xử lý và nghiệp vụ}

\section{Ngôn ngữ truy vấn và suy diễn đặc thù}

\subsection{An toàn truy vấn}

An toàn có nhiều tầng. Pydantic giới hạn request length/type; parser chỉ nhận
chín intent; mỗi intent ánh xạ đến một Cypher pattern do ứng dụng sở hữu; entity
value nằm trong Neo4j parameter; public API không có write operation. Không có
dữ liệu thô nào do người dùng nhập được dùng làm cấu trúc Cypher.

Parameterized query ngăn injection qua value nhưng không tự ngăn query đắt. Vì
vậy query catalog còn áp giới hạn kết quả, shortest-path depth~$\leq 8$,
fixed pattern và không
cho arbitrary Cypher. Query safety không thể chỉ mô tả bằng một câu “đã dùng
parameter”.

\section{Bộ truy vấn và thao tác nghiệp vụ từ cơ bản đến nâng cao}

\subsection{Cypher catalog và CRUD}

Catalog có 10 query từ lookup đến shortest path, aggregation và kiểm tra fact suy
ra. CRUD Movie gồm create/upsert, read neighborhood, update metadata, guarded
delete và upsert \texttt{ACTED\_IN}; mọi value là parameter. CRUD không mở qua public API.

\subsection{Competency questions}

\begin{enumerate}
\item Phim nào do một đạo diễn thực hiện?
\item Ai đóng trong một phim?
\item Hai diễn viên có phim chung nào?
\item Phim thuộc thể loại X có rating trên ngưỡng?
\item Ai từng đóng chung với một người?
\item Đạo diễn nào thường làm phim thuộc thể loại X?
\item Đường ngắn nhất giữa hai người?
\item Phim nào tương tự một phim nguồn và vì sao?
\item Thống kê số node theo label?
\item Fact co-star nào được suy ra và bằng chứng là phim nào?
\end{enumerate}

\subsection{Ma trận truy vết competency question}

\begin{longtable}{|>{\raggedright\arraybackslash}p{0.22\textwidth}|>{\raggedright\arraybackslash}p{0.22\textwidth}|>{\raggedright\arraybackslash}p{0.22\textwidth}|>{\raggedright\arraybackslash}p{0.22\textwidth}|}
\caption{Tổng hợp ma trận truy vết competency question}\label{tab:4-1} \\
\hline
\textbf{CQ} & \textbf{Pattern chính} & \textbf{Số bước} & \textbf{Đầu ra/bằng chứng} \\ \hline
\endfirsthead\hline \textbf{CQ} & \textbf{Pattern chính} & \textbf{Số bước} & \textbf{Đầu ra/bằng chứng} \\ \hline\endhead
1 & Person-DIRECTED-Movie & 1 & Movie ID/title, director link \\ \hline
2 & Person-ACTED\_IN-Movie & 1 & Person và edge metadata \\ \hline
3 & Person-ACTED\_IN-Movie-ACTED\_IN-Person & 2 & shared Movie \\ \hline
4 & Movie-HAS\_GENRE-Genre + rating & 1 + filter & Movie/rating/genre \\ \hline
5 & Person-CO\_STARRED\_WITH-Person & derived & count/evidence IDs \\ \hline
6 & Person-DIRECTED-Movie-HAS\_GENRE-Genre & 2 + aggregate & director/count \\ \hline
7 & variable path giữa Person & tối đa 8 edge & node/relationship path \\ \hline
8 & Movie-feature-Movie & 2 mỗi feature & contribution/feature \\ \hline
9 & scan label & 0 & count theo label \\ \hline
10 & CO\_STARRED\_WITH & derived audit & pair/count/evidence IDs \\ \hline
\end{longtable}

Ma trận là cầu nối giữa yêu cầu, schema, query catalog, API và test. Class hoặc
relationship không hỗ trợ competency question nào cần được cân nhắc loại khỏi
MVP; một CQ không ánh xạ được query chỉ ra thiếu sót thiết kế.

\subsection{Acceptance criteria}

Dataset phải đạt ít nhất 2.000 Movie; import hai lần không đổi counts; validation
không có violation chưa giải thích; catalog Cypher phải parse/chạy; API trả
schema hợp lệ; recommendation trả explanation; evaluation ghi backend, corpus và
quy trình thực nghiệm; kiểm thử tích hợp phải kết nối một Neo4j kiểm thử riêng
thay vì bỏ qua do đồ thị
demo có dữ liệu.

\section{Truy vấn có suy luận và tri thức ẩn}

\subsection{Suy diễn \texttt{CO\_STARRED\_WITH}}
\reportfigure{costar_reasoning.pdf}{Suy diễn quan hệ đồng diễn và đường bằng chứng}{fig:costar-reasoning}

Luật:

\begin{Verbatim}[fontsize=\scriptsize,frame=single]
Person A -[:ACTED_IN]-> Movie <-[:ACTED_IN]- Person B
=> A -[:CO_STARRED_WITH]-> B
\end{Verbatim}

Điều kiện sắp thứ tự ID tránh cặp trùng/ngược. Relationship lưu \texttt{movie\_count},
\texttt{evidence\_movie\_ids} và \texttt{derived=true}. \texttt{MERGE} bảo đảm materialization idempotent.
Đây là rule nghiệp vụ được vật chất hóa trực tiếp bằng Cypher.
